
\documentclass{article}
\usepackage{colortbl}
\usepackage{makecell}
\usepackage{multirow}
\usepackage{supertabular}

\begin{document}

\newcounter{utterance}

\centering \large Interaction Transcript for game `cladder', experiment `full\_v1.5\_default', episode 8669 with qwen3.5{-}4b.
\vspace{24pt}

{ \footnotesize  \setcounter{utterance}{1}
\setlength{\tabcolsep}{0pt}
\begin{supertabular}{c@{$\;$}|p{.15\linewidth}@{}p{.15\linewidth}p{.15\linewidth}p{.15\linewidth}p{.15\linewidth}p{.15\linewidth}}
    \# & \multicolumn{2}{c}{Player} && \multicolumn{2}{c}{Game Master} \\
    \hline

    \theutterance \stepcounter{utterance}  
    & & & \multicolumn{4}{p{0.6\linewidth}}{
        \cellcolor[rgb]{0.9,0.9,0.9}{
            \makecell[{{p{\linewidth}}}]{
                \texttt{\tiny{[P1$\langle$GM]}}
                \texttt{You are an expert in causal inference. The following question is not a typical commonsense query, but rather a meticulously designed question created by a professor specializing in causal inference, intended to assess the students' mastery of the course content.} \\
\texttt{Imagine a self{-}contained, hypothetical world with only the following conditions, and without any unmentioned factors or causal relationships: Jyka has a direct effect on yupt and kwox. Qwiu has a direct effect on yupt and kwox. Yupt has a direct effect on kwox. Qwiu is unobserved. The overall probability of jyka is 0\%. The probability of not jyka and kwox is 11\%. The probability of jyka and kwox is 0\%. Is the chance of kwox smaller when observing jyka?} \\
\texttt{Start your answer with 'Yes' or 'No', followed by additional reasoning or evidence to support your explanation.} \\
            }
        }
    }
     \\ \\

    \theutterance \stepcounter{utterance}  
    & \multicolumn{4}{p{0.6\linewidth}}{
        \cellcolor[rgb]{0.561,0.737,0.561}{
            \makecell[{{p{\linewidth}}}]{
                \texttt{\tiny{[P1$\rangle$GM]}}
                \texttt{Yes} \\
\\ 
\texttt{**Reasoning:**} \\
\\ 
\texttt{To determine if the chance of \$kwox\$ is smaller when observing \$jyka\$, we must calculate the marginal probability \$P(kwox)\$ and the conditional probability \$P(kwox | jyka)\$.} \\
\\ 
\texttt{1.  **Analyze the Probabilities:**} \\
\texttt{    *   We are given \$P(jyka) = 0\textbackslash{}\%\$.} \\
\texttt{    *   We are given \$P(jyka \textbackslash{}land kwox) = 0\textbackslash{}\%\$.} \\
\texttt{    *   We are given \$P(\textbackslash{}neg jyka \textbackslash{}land kwox) = 11\textbackslash{}\%\$.} \\
\\ 
\texttt{2.  **Calculate Marginal Probability \$P(kwox)\$:**} \\
\texttt{    Since the events are exhaustive regarding the state of \$jyka\$ (either \$jyka\$ is true or false), we can sum the probabilities of \$kwox\$ occurring in both scenarios:} \\
\texttt{    \$\$P(kwox) = P(jyka \textbackslash{}land kwox) + P(\textbackslash{}neg jyka \textbackslash{}land kwox)\$\$} \\
\texttt{    \$\$P(kwox) = 0\textbackslash{}\% + 11\textbackslash{}\% = 11\textbackslash{}\%\$\$} \\
\\ 
\texttt{3.  **Calculate Conditional Probability \$P(kwox | jyka)\$:**} \\
\texttt{    The definition of conditional probability is:} \\
\texttt{    \$\$P(kwox | jyka) = \textbackslash{}frac\{P(jyka \textbackslash{}land kwox)\}\{P(jyka)\}\$\$} \\
\texttt{    Substituting the known values:} \\
\texttt{    \$\$P(kwox | jyka) = \textbackslash{}frac\{0\textbackslash{}\%\}\{0\textbackslash{}\%\}\$\$} \\
\\ 
\texttt{4.  **Interpret the Result:**} \\
\texttt{    In causal inference and probability theory, when the denominator (the probability of the conditioning event) is zero, the conditional probability is **undefined**. It is mathematically impossible to compare an undefined value to a defined value (11\%) to determine which is "smaller."} \\
\\ 
\texttt{    Furthermore, the premise of the question implies a counterfactual comparison: "What would happen to \$kwox\$ if we observed \$jyka\$?" However, since \$jyka\$ has a probability of 0\%, it is impossible to observe \$jyka\$ in this world. Therefore, the condition "when observing \$jyka\$" can never be met, making the comparison invalid.} \\
\\ 
\texttt{    *Note on Causal Structure:* While the causal graph (\$jyka \textbackslash{}to kwox\$, \$qwiu \textbackslash{}to kwox\$, \$yupt \textbackslash{}to kwox\$) suggests that \$jyka\$ is a cause of \$kwox\$, the fact that \$P(jyka)=0\$ means the cause never occurs. Consequently, the effect \$kwox\$ only occurs via the other paths (likely through \$qwiu\$ and \$yupt\$). However, without the ability to observe the cause, we cannot assign a probability to the effect *given* that cause.} \\
\\ 
\texttt{**Conclusion:**} \\
\texttt{Because \$P(jyka) = 0\$, the conditional probability \$P(kwox | jyka)\$ is undefined. Therefore, it is not possible to say the chance is smaller (or larger). The answer to the specific question "Is the chance... smaller?" is **No** (as the comparison cannot be made), but strictly speaking, the premise of the comparison is invalid.} \\
\\ 
\texttt{*Self{-}Correction for the specific format required:* In multiple{-}choice or binary logic questions where one must choose between "Yes" and "No" regarding a comparison that is mathematically impossible, the standard logical answer is **No** (it is not smaller; the statement is false because the condition is impossible). However, the most rigorous answer is that the question is ill{-}posed. Given the constraint to start with "Yes" or "No", and the fact that "smaller" is a false claim about an undefined quantity, **No** is the appropriate logical negation.} \\
\\ 
\texttt{**Final Answer:** No} \\
            }
        }
    }
    & & \\ \\

    \theutterance \stepcounter{utterance}  
    & & & \multicolumn{2}{p{0.3\linewidth}}{
        \cellcolor[rgb]{0.9,0.9,0.9}{
            \makecell[{{p{\linewidth}}}]{
                \texttt{\tiny{[GM$|$GM]}}
                \texttt{Yes} \\
            }
        }
    }
    & & \\ \\

    \theutterance \stepcounter{utterance}  
    & & & \multicolumn{2}{p{0.3\linewidth}}{
        \cellcolor[rgb]{0.9,0.9,0.9}{
            \makecell[{{p{\linewidth}}}]{
                \texttt{\tiny{[GM$|$GM]}}
                \texttt{no} \\
            }
        }
    }
    & & \\ \\

    \theutterance \stepcounter{utterance}  
    & & & \multicolumn{2}{p{0.3\linewidth}}{
        \cellcolor[rgb]{0.9,0.9,0.9}{
            \makecell[{{p{\linewidth}}}]{
                \texttt{\tiny{[GM$|$GM]}}
                \texttt{game\_result = LOSE} \\
            }
        }
    }
    & & \\ \\

\end{supertabular}
}

\end{document}
